\subsubsection*{定理 2.2.}
\underline{设 $X$ 在域 $k$ 上光滑, $\Delta$ 为余维 $1+r$、处处光滑的子概形, $f:\widetilde X\longrightarrow X$ 为 $X$ 沿 $\Delta$ 的爆破. 则:}

\medskip
\noindent(2.2.1) \underline{Leray 谱序列}
\[
E_2^{p,q}=H^p\bigl(X,R^qf_*\mathbb Z_\ell\bigr)
\Longrightarrow
H^{p+q}\bigl(\widetilde X,\mathbb Z_\ell\bigr)
\]
\underline{退化.}

\medskip
\noindent(2.2.2) \underline{存在加法同构}
\[
\begin{aligned}
H^q(\widetilde X,\mathbb Z_\ell)
&\xleftarrow{\ \sim\ }
H^q(X,\mathbb Z_\ell)\oplus
\frac{H^q(\widetilde\Delta,\mathbb Z_\ell)}
     {\operatorname{Im}H^q(\Delta,\mathbb Z_\ell)}
\\
&\simeq
\left(\bigoplus_{j=1}^{r}
H^{q-2j}(\Delta,\mathbb Z_\ell(-j))\right)
\oplus H^q(X,\mathbb Z_\ell).
\end{aligned}
\]

\noindent\underline{证明.}
设
\begin{equation}\tag{2.2.3}
g:\widetilde\Delta\longrightarrow\Delta
\end{equation}
为 $f$ 在 $\widetilde\Delta$ 上的限制, 并设
\begin{equation}\tag{2.2.4}
\left\{
\begin{aligned}
j&:\Delta\hookrightarrow X,\\
i&:\widetilde\Delta\hookrightarrow\widetilde X
\end{aligned}
\right.
\end{equation}
为包含映射, 从而有交换图
\begin{equation}\tag{2.2.5}
\begin{tikzcd}[column sep=huge,row sep=large]
\widetilde X \arrow[d,"f"'] &
\widetilde\Delta \arrow[l,"i"'] \arrow[d,"g"]\\
X & \Delta \arrow[l,"j"']
\end{tikzcd}.
\end{equation}
它在态射 $f$ 和 $g$ 关于 $\widetilde X$ 和 $\widetilde\Delta$ 上的层 $\mathbb Z_\ell$ 的 Leray 谱序列之间诱导同态:
\begin{equation}\tag{2.2.6}
E_s^{p,q}(f)\xrightarrow{\ (i^*,j^*)\ }E_s^{p,q}(g).
\end{equation}
另一方面, 由于 $f$ 诱导 $\widetilde X-\widetilde\Delta$ 到 $X-\Delta$ 的同构, 固有基变换定理给出
\begin{equation}\tag{2.2.7}
R^qf_*\mathbb Z_\ell\simeq
\begin{cases}
\mathbb Z_\ell,&q=0,\\[1ex]
j_*R^qg_*\mathbb Z_\ell,&q>0.
\end{cases}
\end{equation}
因此
\begin{equation}\tag{2.2.8}
E_2^{p,q}(f)
\xrightarrow[\ \sim\ ]{\ (i^*,j^*)\ }
E_2^{p,q}(g)
\qquad(q\ne0).
\end{equation}
我们将对 $s$ 归纳证明
\begin{equation}\tag{2.2.9}
E_s^{p,q}(f)
\xrightarrow[\ \sim\ ]{\ (i^*,j^*)\ }
E_s^{p,q}(g)
\qquad(q\ne0,\ s\geq2)
\end{equation}
以及
\begin{equation}\tag{2.2.10}
d_s^{p,q}(f)=0\qquad(s\geq2).
\end{equation}
暂且承认公式
\begin{equation}\tag{2.2.11}
d_s^{p,s-1}(f)=0\qquad(s\geq2,\ \text{对每个 }p).
\end{equation}

先考虑 $s=2$. 有交换图
\begin{equation}\tag{2.2.12}
\begin{tikzcd}[column sep=huge,row sep=large]
E_2^{p,q}(f)
  \arrow[r,"{(i^*,j^*)}"]
  \arrow[d,"{d_2^{p,q}(f)}"]
&
E_2^{p,q}(g)
  \arrow[d,"{0=d_2^{p,q}(g)}"]
\\
E_2^{p+2,q-1}(f)
  \arrow[r,"{(i^*,j^*)}"]
&
E_2^{p+2,q-1}(g)
\end{tikzcd}.
\end{equation}
当 $q\leq0$ 时, 有 $E_2^{p+2,q-1}(f)=0$, 因而 $d_2^{p,q}(f)=0$. 当 $q=1$ 时, 公式 (2.2.11) 给出 $d_2^{p,1}(f)=0$. 当 $q>1$ 时, (2.2.12) 中的水平箭头为同构, 而由 1.2 和 2.1.4 有 $d_2^{p,q}(g)=0$. 这就证明了 $s=2$ 时的 (2.2.10), 而 $s=2$ 时的 (2.2.9) 恰为 (2.2.8).

假设对某个 $s\geq2$ 已知 (2.2.9) 和 (2.2.10). 假设 (2.2.10) 蕴涵 $E_s^{p,q}(f)\simeq E_{s+1}^{p,q}(f)$, 又由 1.2 有 $E_s^{p,q}(g)\simeq E_{s+1}^{p,q}(g)$, 所以有交换图
\begin{equation}\tag{2.2.13}
\begin{tikzcd}[column sep=huge,row sep=large]
E_s^{p,q}(f)
  \arrow[r,"{(i^*,j^*)}"]
  \arrow[d,"\wr"']
&
E_s^{p,q}(g)
  \arrow[d,"\wr"]
\\
E_{s+1}^{p,q}(f)
  \arrow[r,"{(i^*,j^*)}"]
&
E_{s+1}^{p,q}(g)
\end{tikzcd},
\end{equation}
结合 $s$ 时的 (2.2.9), 这就证明了 $s+1$ 时的 (2.2.9). 为推出 $s+1$ 时的 (2.2.10), 考虑交换图
\begin{equation}\tag{2.2.14}
\begin{tikzcd}[column sep=huge,row sep=large]
E_{s+1}^{p,q}(f)
  \arrow[r,"{(i^*,j^*)}"]
  \arrow[d,"{d_{s+1}^{p,q}(f)}"]
&
E_{s+1}^{p,q}(g)
  \arrow[d,"{0=d_{s+1}^{p,q}(g)}"]
\\
E_{s+1}^{p+s+1,q-s}(f)
  \arrow[r,"{(i^*,j^*)}"]
&
E_{s+1}^{p+s+1,q-s}(g)
\end{tikzcd}.
\end{equation}
当 $q\leq s-1$ 时, 有 $E_{s+1}^{p+s+1,q-s}(f)=0$, 因而 $d_{s+1}^{p,q}(f)=0$. 当 $q=s$ 时, 有 (2.2.11); 而当 $q>s$ 时, (2.2.14) 中的水平箭头为同构, 且由 1.2 有 $d_{s+1}^{p,q}(g)=0$. 这就得到所需的结论.
\footnote{印本在此处写作 $d_s^{p,q}(f)=0$, 但归纳步骤和图 (2.2.14) 涉及的是 $d_{s+1}^{p,q}(f)$.}

尚须验证 (2.2.11), 它的意思是复合边缘同态
\begin{equation}\tag{2.2.15}
E_2^{p,0}(f)\longrightarrow E_3^{p,0}(f)
\longrightarrow E_4^{p,0}(f)\longrightarrow\cdots
\longrightarrow E_\infty^{p,0}(f)
\end{equation}
为同构; 即
\begin{equation}\tag{2.2.16}
f^*:H^p(X,\mathbb Z_\ell)
\longrightarrow H^p(\widetilde X,\mathbb Z_\ell)
\end{equation}
为单射. 但它有一个左逆, 即\underline{Gysin 同态} (SGA 5 IV)
\begin{equation}\tag{2.2.17}
f_*:H^p(\widetilde X,\mathbb Z_\ell)
\longrightarrow H^p(X,\mathbb Z_\ell),
\end{equation}
它确实满足
\begin{equation}\tag{2.2.18}
f_*f^*=\mathrm{id}.
\end{equation}
这由投影公式和 $f_*(1)=1$ 这一事实得出, 因为 $f$ 为双有理态射, 所以其次数为一. 这就证明了 (2.2.1).

